% CVT Foundations v0.2 — proposed Formal Architecture replacement
% Review artifact. This section replaces only the material between
% \section{Formal Architecture} and \section{Membrane Formulation}.

\section{Formal Architecture: Observables, Gates, and Candidate Models}
\label{sec:formal}

The formal architecture must separate what is measured from how it is normalized and from how normalized quantities are combined. CVT therefore distinguishes three layers:
\begin{enumerate}[label=(\roman*)]
    \item \textbf{observable layer}: domain-native measurements with units, time windows, and uncertainty;
    \item \textbf{gate layer}: documented transformations of those measurements into normalized scores;
    \item \textbf{model layer}: candidate aggregation and dynamical rules whose performance can be compared against an independently defined viability outcome.
\end{enumerate}
A normalized gate is not itself an observation, and no scalar aggregator defines success merely by being chosen.

\subsection{Observable and Gate Layers}

For each proposed condition \(i\in\mathcal I=\{B,Q,C,S\}\), let \(z_i(t)\) denote a raw domain-specific measurement or vector of measurements. A normalization map
\begin{equation}
    G_i(t)
    =
    g_i\!\left(z_i(t);\vartheta_i\right)
    \in[0,1]
    \label{eq:gate-normalization}
\end{equation}
converts the raw measurement into a gate score. The parameters \(\vartheta_i\), normalization range, uncertainty interval, and threshold rationale must be reported for each application.

The maps need not be monotone. In particular, bounded coupling generally requires a two-sided map: both insufficient exchange and excessive exchange may receive low scores, while an admissible operating band receives a high score. Preserved organization \(Q\), capture readiness \(C\), and available restorative reserve \(S\) likewise require domain-native observables rather than intuitive numerical assignments.

Let \(y(t)\) denote a domain-native viability outcome and let \(\mathcal Y_{\mathrm{viable}}\) be declared before the gates are fitted. The binary event
\begin{equation}
    V_{\mathrm{set}}(t)
    =
    \mathbf{1}\!\left\{y(t)\in\mathcal Y_{\mathrm{viable}}\right\}
    \label{eq:independent-viability-event}
\end{equation}
is an independently defined outcome, not a CVT aggregate. Continuous domain-native outcomes may be used instead when binary classification would discard relevant information.

\subsection{Candidate Aggregation Family}

The intersection in Equation~\eqref{eq:candidate-cvt-region} specifies simultaneous constraint satisfaction. When a scalar score is useful, several non-compensatory or weakly compensatory representations should be compared:
\begin{align}
    \V_{\Pi}(t)
    &=
    \prod_{i\in\mathcal I}G_i(t),
    \\
    \V_{\min}(t)
    &=
    \min_{i\in\mathcal I}G_i(t),
    \\
    \V_G(t)
    &=
    \prod_{i\in\mathcal I}G_i(t)^{w_i},
    \qquad
    w_i\geq0,
    \quad
    \sum_iw_i=1,
    \\
    \V_{\tau}(t)
    &=
    -\tau\log\!\left(
        \frac{1}{|\mathcal I|}
        \sum_{i\in\mathcal I}e^{-G_i(t)/\tau}
    \right),
    \qquad \tau>0.
    \label{eq:aggregation-family}
\end{align}
The last expression is a normalized soft minimum. Direct simultaneous constraints \(G_i\geq\Theta_i\) for every essential condition remain a further alternative. Product, minimum, weighted geometric mean, soft minimum, and constraint intersection encode different assumptions and must be evaluated against the same independently defined outcome using common data or simulations and out-of-sample comparison.

\subsection{Threshold Violations}

Given a gate threshold \(\Theta_i\), define the one-sided violation
\begin{equation}
    \Phi_i(t)
    =
    [\Theta_i-G_i(t)]_+,
    \qquad
    [x]_+=\max\{x,0\}.
    \label{eq:violation-variable}
\end{equation}
A two-sided raw constraint should be encoded in the normalization map \(g_i\), so that \(\Phi_i\) retains one interpretation: distance below the normalized viability threshold. Thresholds should be justified independently, preregistered when possible, or estimated out of sample. The uncertainty in \(\Phi_i\) inherits uncertainty in both \(z_i\) and \(\Theta_i\).

\subsection{Candidate Gate Dynamics}

A phenomenological recovery model may be written as
\begin{equation}
    \dot G_i
    =
    a_i(H,S)\Phi_i
    -
    \ell_i(t)
    -
    \sum_{j\neq i}\kappa_{ij}\Phi_j,
    \label{eq:gate-dynamics}
\end{equation}
with projected or saturating dynamics used to preserve \(G_i\in[0,1]\). Here \(\ell_i(t)\) is a normalized gate-degradation rate with units of inverse time, not an unscaled physical load. The coefficients \(\kappa_{ij}\) represent lagged cross-gate degradation and cannot be inferred from simultaneous decline alone, because common-cause forcing is an alternative explanation.

One candidate recovery coefficient is
\begin{equation}
    a_i(H,S)
    =
    \alpha_i(1+\eta_i H)q_i(S),
    \label{eq:recovery-coefficient}
\end{equation}
where \(q_i\) is nondecreasing and may satisfy \(q_i(0)=0\) when recovery requires available reserve. Equation~\eqref{eq:gate-dynamics} is not a universal law. It is a model class whose terms must be mapped to domain-native mechanisms and compared with simpler baselines.

\subsection{Recovery History and Available Reserve}

Recovery history or maturity and present restorative reserve are distinct states. Let
\[
    H(t)\in[0,1]
\]
denote learned recovery history, conditioning, or maturity, and let
\[
    S(t)\in[0,S_{\max}]
\]
denote currently available restorative reserve. The normalized reserve gate is \(G_S=g_S(S)\).

A candidate learning law is
\begin{equation}
    \dot H
    =
    \epsilon U_{\mathrm{rec}}(t)(1-H)
    -
    \zeta H
    -
    \xi_D D H,
    \label{eq:history-dynamics}
\end{equation}
where \(U_{\mathrm{rec}}(t)\) is an independently coded measure of successful viable operation or recovery. A candidate reserve law is
\begin{equation}
    \dot S
    =
    r_S(H)(S_{\max}-S)
    -
    c_S(t)
    -
    \sum_i\nu_i\Phi_i,
    \label{eq:reserve-dynamics}
\end{equation}
with projection onto \([0,S_{\max}]\). The term \(c_S(t)\) is measured reserve consumption; \(r_S(H)\) is a replenishment rate that may improve with history. These equations permit a mature system to be temporarily exhausted and an inexperienced system to possess substantial unused reserve. Neither possibility was representable when memory, maturity, and reserve shared one variable.

\subsection{Raw Exchange, Productive Uptake, and Prediction}

Let \(J_{\mathrm{raw}}(t)\) be measured delivered exchange. In a transport model one may have
\begin{equation}
    J_{\mathrm{raw}}
    =
    \Pi\Delta,
    \label{eq:raw-exchange}
\end{equation}
where \(\Pi\) is permeability or exchange capacity and \(\Delta\) is a driving gradient. This relation is domain-specific rather than the definition of raw exchange in every application.

Let \(J_{\mathrm{eff}}(t)\) be independently measured retained productive uptake over the same accounting window. When the quantities are commensurate and \(J_{\mathrm{raw}}>0\), observed uptake efficiency is
\begin{equation}
    \rho_{\mathrm{obs}}(t)
    =
    \frac{J_{\mathrm{eff}}(t)}{J_{\mathrm{raw}}(t)}.
    \label{eq:effective-exchange}
\end{equation}
In systems with endogenous amplification or noncommensurate input and output units, productive uptake must instead be defined through an explicit conversion or target-response measure.

Capture readiness is measured before or independently of the uptake outcome. CVT may then test a predictor
\begin{equation}
    \widehat{\rho}_{\theta}(t)
    =
    f_{\theta}\!\left(
        G_B(t),G_Q(t),G_C(t),G_S(t),\mathbf w(t)
    \right),
    \label{eq:uptake-predictor}
\end{equation}
where \(\mathbf w(t)\) contains declared exogenous covariates. The predictor may be multiplicative, minimum-gated, nonlinear, or domain-native. Defining \(\rho\) as a product of the same gates and then treating that product as an explanation of uptake is not permitted.

\subsection{Damage Accumulation}

Reserve the symbol \(D(t)\in[0,1]\) for accumulated damage. When raw and effective exchange are commensurate, define unretained delivered load as
\begin{equation}
    L_{\mathrm{u}}(t)
    =
    [J_{\mathrm{raw}}(t)-J_{\mathrm{eff}}(t)]_+.
    \label{eq:unretained-load}
\end{equation}
Otherwise \(L_{\mathrm{u}}\) must be replaced by a directly measured rejection, spillover, waste, or harmful-load quantity.

A candidate damage model is
\begin{equation}
    \dot D
    =
    \gamma L_{\mathrm{u}}
    +
    \delta\sum_i\Phi_i
    +
    \omega R_{\mathrm{risk}}
    -
    \mu(H,S)D,
    \label{eq:damage}
\end{equation}
with projected dynamics preserving \(D\in[0,1]\). The domain-specific risk term \(R_{\mathrm{risk}}\) must be specified before fitting rather than used as a residual container. The repair rate \(\mu(H,S)\) may depend on both learned recovery capacity and currently available reserve. Equation~\eqref{eq:damage} is a candidate mechanism, not a general conservation law.

\subsection{Hosted Transformation as an Independent Outcome}

\begin{definition}[Hosted transformation]
Fix a host state \(x(t)\), a transformation variable \(X(t)\), a domain-native viable set \(\mathcal K_{\mathrm{declared}}\), a transition threshold \(X_{\mathrm{crit}}\), and a follow-up horizon \(\tau_f\geq0\). A transformation is hosted on \([t_0,t_1]\) when
\begin{equation}
    X(t_1)-X(t_0)
    \geq
    X_{\mathrm{crit}},
    \label{eq:transformation-success}
\end{equation}
and
\begin{equation}
    x(t)\in\mathcal K_{\mathrm{declared}}
    \qquad
    \text{for all }t\in[t_0,t_1+\tau_f].
    \label{eq:delayed-host-viability}
\end{equation}
An independently specified damage limit may be added when damage is not already included in \(\mathcal K_{\mathrm{declared}}\).
\end{definition}

The follow-up horizon permits delayed failure to count against apparent success. Gate scores and aggregate CVT scores are predictors or explanatory variables for this outcome; they do not define the outcome they are supposed to predict.

\subsection{Host-Paced Transition as a Model Class}

Let \(u_X(t)\) denote the control or pacing input applied to a transformation. A candidate host-paced policy is
\begin{equation}
    u_X(t)
    =
    u_{\max}
    h_{\theta}\!\left(
        G_B,G_Q,G_C,G_S,D
    \right),
    \qquad
    0\leq h_{\theta}\leq1,
    \label{eq:host-capacity-policy}
\end{equation}
combined with domain-specific transition dynamics
\begin{equation}
    \dot X
    =
    F\!\left(X,u_X,\mathbf w\right).
    \label{eq:host-gated-transition}
\end{equation}
The pacing function \(h_{\theta}\) must not double-count variables already embedded in one another. It should be compared with fixed-rate, capacity-maximizing, and domain-native control policies using the same outcome criteria. CVT does not imply the dimensionally unsupported universal inequality \(\dot X\leq R_{\mathrm{host}}\); any rate bound must be derived in the units and dynamics of the specific application.

\subsection{Minimum Viable Host as a Domain-Specific Threshold}

For a fixed model, forcing severity \(\Lambda\), observation horizon \(\tau\), and control policy \(\pi\), one may estimate a reserve threshold
\begin{equation}
    S(t_0)
    \geq
    S_{\mathrm{crit}}(\Lambda,\tau;\pi,\theta)
    \label{eq:minimum-host-threshold}
\end{equation}
associated with a chosen probability or level of hosted success. This is a domain- and model-specific threshold, not a general theorem. Its form may be nonlinear, history-dependent, or absent. A case that hosts the transformation with no measurable restorative reserve would constrain the claim that reserve is essential in that domain.

\subsection{Local Exchange--Uptake Condition}

Suppose \(J_{\mathrm{raw}}>0\), \(\rho_{\mathrm{obs}}>0\), and
\[
    J_{\mathrm{eff}}
    =
    \rho_{\mathrm{obs}}(J_{\mathrm{raw}})J_{\mathrm{raw}}
\]
locally with differentiable \(\rho_{\mathrm{obs}}\). Then
\begin{equation}
    \frac{dJ_{\mathrm{eff}}}{dJ_{\mathrm{raw}}}
    =
    \rho_{\mathrm{obs}}
    +
    J_{\mathrm{raw}}
    \frac{d\rho_{\mathrm{obs}}}{dJ_{\mathrm{raw}}}.
    \label{eq:exchange-uptake-derivative}
\end{equation}
Effective uptake decreases with additional delivered exchange exactly when
\begin{equation}
    \frac{d\log\rho_{\mathrm{obs}}}{d\log J_{\mathrm{raw}}}
    <
    -1.
    \label{eq:uptake-elasticity-condition}
\end{equation}
This is a calculus condition, not evidence that such degradation occurs. The empirical CVT claim is that some scoped systems enter regimes satisfying Equation~\eqref{eq:uptake-elasticity-condition} because capture readiness, restorative reserve, or another essential mechanism deteriorates under load.

The formal research burden is therefore explicit: define viability independently, publish the raw measurements and normalization maps, estimate uncertainty, compare aggregation and pacing models out of sample, and search deliberately for transformations that remain hosted despite failure of a proposed essential condition.
